\chapter{Текстура}\label{ch:texture}

Поликристалл состоит из миллионов зёрен. Если их решётки повёрнуты как попало, металл в среднем изотропен. Но
после прокатки и отжига ориентации зёрен распределены неслучайно --- у металла есть \textbf{текстура}. В цирконии,
где каждое зерно сильно анизотропно, текстура определяет и прочность трубы в разных направлениях, и то, как лягут
в ней гидриды. Эта глава --- о том, как текстуру описывают, измеряют и откуда она берётся.

\section{Ориентация зерна}

Для трубы естественная система координат образца --- \textbf{R} (радиус), \textbf{T} (окружное направление) и
\textbf{A} (ось трубы). У каждого зерна своя система координат, связанная с решёткой: для ГПУ удобно взять
$\mathbf e_1\parallel\mathbf a_1$, $\mathbf e_3\parallel\mathbf c$ и $\mathbf e_2$ перпендикулярно им. \textbf{Ориентация}
зерна --- это поворот, совмещающий систему образца с системой кристалла. Её задают матрицей поворота $\mathbf g$
(3$\times$3, ортогональная) или тремя углами.

Чаще всего используют \textbf{углы Эйлера} в записи Бунге $(\varphi_1, \Phi, \varphi_2)$: поворот на $\varphi_1$ вокруг оси
$Z$ образца, затем на $\Phi$ вокруг новой оси $X$, затем на $\varphi_2$ вокруг новой оси $Z$. Ось $c$ кристалла в
координатах образца тогда
\begin{equation}
\mathbf c = (\sin\varphi_1\sin\Phi,\; -\cos\varphi_1\sin\Phi,\; \cos\Phi),
\end{equation}
то есть $\Phi$ --- угол между осью $c$ и осью $Z$ образца, а $\varphi_2$ --- поворот кристалла вокруг собственной
оси $c$. Для ГПУ из-за шестикратной симметрии $\varphi_2$ достаточно брать от 0 до 60\textdegree.

\fig{t1_single}{а) Одно зерно (ГПУ призма) в осях трубы R, T, A; ось $c$ наклонена от радиуса к окружному
направлению на 35\textdegree. б) Его полюсы на полюсной фигуре (центр --- R, по горизонтали --- T, по вертикали --- A):
один полюс $(0002)$, три полюса $\{10\bar10\}$ и три $\{11\bar20\}$ в видимой полусфере. в) Полюсы $(0002)$ 400 зёрен
трубы.}

\section{Полюсные фигуры}

Нарисовать ориентацию зерна целиком трудно, но легко нарисовать направление нормали к одной выбранной плоскости
--- её полюс на стереографической (глава~\ref{ch:lattice}) или равноплощадной проекции. \textbf{Полюсная фигура}
$\{hkil\}$ --- распределение полюсов этого семейства от всех зёрен в осях образца (рис.~\ref{fig:t1_single}).

Важная тонкость: одно зерно даёт столько полюсов, сколько в семействе эквивалентных плоскостей. Для $(0002)$ ---
один (с противоположным --- два), для $\{10\bar10\}$ --- три пары, для $\{10\bar11\}$ --- шесть пар. Поэтому полюсная фигура
призм выглядит «размазанной», даже когда оси $c$ собраны в узкие пятна.

Когда зёрен тысячи, рисуют не точки, а плотность полюсов в \textbf{долях от случайной} (m.r.d., multiples of random
distribution): 1 --- столько же, сколько при беспорядочной ориентации, 5 --- в пять раз больше. Для этого и нужна
равноплощадная проекция: в ней одинаковые телесные углы занимают одинаковую площадь.

\fig{t2_polefigs}{Полюсные фигуры схематичной текстуры трубы (равноплощадная проекция, центр --- R, справа --- T,
вверху --- A; уровни --- в долях случайной; тонкая линия --- уровень 1). Оси $c$ наклонены от R к T на
$\pm35^\circ$ (разброс $12^\circ$), к оси трубы --- на $15^\circ$. а) $(0002)$: два максимума по обе стороны от центра.
б), в) $\{10\bar10\}$ и $\{11\bar20\}$ холоднодеформированной трубы: направление $\langle10\bar10\rangle$ вдоль оси трубы.
г) $\{10\bar10\}$ после рекристаллизации: кристаллы повёрнуты на 30\textdegree{} вокруг $c$, вдоль оси трубы теперь
$\langle11\bar20\rangle$.}

Для оболочечных труб из циркониевых сплавов характерна картина рис.~\ref{fig:t2_polefigs}. Базисная полюсная
фигура $(0002)$ имеет два максимума, отклонённых от радиального направления к окружному на 30--40\textdegree;
вдоль оси трубы полюсов почти нет. Это значит: оси $c$ лежат около плоскости R--T, а базисные плоскости
большинства зёрен почти параллельны поверхности трубы. Распределение вокруг оси $c$ зависит от обработки: после
холодной прокатки вдоль оси трубы выстраивается $\langle10\bar10\rangle$, после рекристаллизации зёрна в среднем
поворачиваются на 30\textdegree{} вокруг $c$, и вдоль оси оказывается $\langle11\bar20\rangle$. Базисная фигура
при этом почти не меняется.

\section{Обратные полюсные фигуры}

Можно задать вопрос наоборот: какое кристаллическое направление стоит вдоль данного направления образца?
\textbf{Обратная полюсная фигура} для направления R --- распределение этого направления в координатах кристалла. Из-за
симметрии решётки всю сферу можно свернуть в маленький \textbf{стандартный треугольник}: для ГПУ это сектор между
$[0001]$, $\langle2\bar1\bar10\rangle$ и $\langle10\bar10\rangle$ (рис.~\ref{fig:t3_ipf}).

\fig{t3_ipf}{Обратные полюсные фигуры той же текстуры для направлений R, T и A (стандартный треугольник ГПУ,
равноплощадная проекция, уровни в долях случайной). Вдоль радиуса стоят направления в 30--40\textdegree{} от $c$;
вдоль окружности --- направления, близкие к базисной плоскости; вдоль оси трубы --- $\langle10\bar10\rangle$.}

Обратные полюсные фигуры удобны, когда важно одно направление образца: направление нагрузки при растяжении, нормаль к
поверхности при рентгеновской съёмке. По рис.~\ref{fig:t3_ipf}б видно, что при окружном растяжении нагрузка
приложена почти в базисной плоскости большинства зёрен --- самом «мягком» положении для призматического скольжения.

\section{Функция распределения ориентаций}

Полная характеристика текстуры --- \textbf{функция распределения ориентаций} (ФРО) $f(\mathbf g)$: доля объёма
зёрен с ориентацией в малой окрестности $\dd\mathbf g$,
\begin{equation}
\frac{\dd V}{V} = f(\mathbf g)\,\dd\mathbf g,\qquad f\equiv1 \text{ для случайной текстуры}.
\end{equation}
Полюсная фигура --- это «проекция» ФРО: для каждой точки $\mathbf y$ фигуры складываются все ориентации, при которых
нормаль $\mathbf h$ к выбранной плоскости попадает в $\mathbf y$. Каждая полюсная фигура теряет информацию о
повороте вокруг $\mathbf h$, но несколько фигур разных плоскостей вместе позволяют восстановить ФРО. Это делают
разложением по сферическим гармоникам (метод Бунге) или прямыми методами; стандартный современный инструмент ---
пакет MTEX. По ФРО считают любые средние: модули упругости, тепловое расширение, коэффициенты Хилла, факторы Кернса.

\section{Как измеряют текстуру}

\textbf{Рентген, метод Шульца.} Дифрактометр настраивают на угол $2\theta$ одного отражения, а образец поворачивают:
наклоняют на угол $\chi$ (0--75\textdegree) и вращают вокруг нормали на угол $\varphi$ (0--360\textdegree). В каждом
положении в отражающее положение попадают зёрна, у которых нормаль к выбранной плоскости направлена по биссектрисе
падающего и отражённого лучей. Интенсивность как функция $(\chi, \varphi)$ и есть полюсная фигура. При больших
наклонах пятно на образце растягивается, интенсивность падает (расфокусировка), и край фигуры за 75--80\textdegree{}
не измеряется --- его восстанавливают по ФРО. Для трубы образец часто собирают из выпрямленных полосок.

\textbf{EBSD} даёт ориентации отдельных зёрен напрямую (глава~\ref{ch:methods}). Для статистически надёжной текстуры
нужно несколько тысяч зёрен, зато видно, \emph{какие} зёрна где стоят, и можно связать ориентацию зерна с тем,
есть ли в нём радиальный гидрид.

\textbf{Нейтроны} проходят сантиметры металла и дают среднюю текстуру по большому объёму.

\textbf{Обратная полюсная фигура по обычной дифрактограмме} (метод Харриса). При съёмке $\theta$--$2\theta$ с плоского
образца каждое отражение приходит только от зёрен, у которых нормаль к плоскости $(hkil)$ совпадает с нормалью к
поверхности. Отношение интенсивности к порошковой даёт плотность полюсов в этом направлении:
\begin{equation}
P_{hkil} = \frac{I_{hkil}/I^0_{hkil}}{\tfrac1N\sum_{j=1}^N I_j/I^0_j},
\end{equation}
где $I^0$ --- интенсивность для порошка, а знаменатель нормирует среднее на единицу. Набор $P_{hkil}$ --- точки на
обратной полюсной фигуре для нормали к поверхности (рис.~\ref{fig:t4_tc}).

\section{Факторы Кернса}

Для гидридов удобна свёртка текстуры в три числа --- \textbf{факторы Кернса}:
\begin{keyf}[Факторы Кернса]
\[
f_d = \bigl\langle\cos^2\angle(\mathbf c,\mathbf d)\bigr\rangle_{\text{по объёму}},\qquad f_R+f_T+f_A = 1 .
\]
\end{keyf}
$f_R$ --- доля «ориентированности» осей $c$ по радиусу, $f_T$ --- по окружности, $f_A$ --- вдоль оси трубы. Сумма равна
единице, потому что квадраты направляющих косинусов любого вектора в сумме дают 1. Для случайной текстуры все три
фактора равны 1/3. Для типичных оболочек $f_R\approx0.5$--$0.7$, $f_T\approx0.25$--$0.35$, $f_A\approx0.05$--$0.1$.

\fig{f7_texture}{Та же схематичная текстура трубы: а) полюсная фигура $(0002)$ (точки --- оси $c$ 4000 зёрен,
центр --- радиальное направление); б) её факторы Кернса.}

Почему именно $\cos^2$? Многие свойства одного зерна --- тепловое расширение, модуль упругости, деформация от
несоответствия гидрида --- являются тензорами второго ранга с осью симметрии $c$. Значение такого свойства вдоль
направления $\mathbf d$ равно $p_\perp + (p_\parallel - p_\perp)\cos^2\angle(\mathbf c,\mathbf d)$, и его среднее по
поликристаллу (в простейшем приближении) выражается ровно через $f_d$. Например, коэффициент теплового расширения
трубы в радиальном направлении $\bar\alpha_R = \alpha_a + (\alpha_c - \alpha_a)f_R$.

Факторы Кернса считают по ФРО или по обратной полюсной фигуре нормали к поверхности: каждому измеренному полюсу
приписывают долю телесного угла стандартного треугольника, ближайшую к нему (ячейки Вороного), и
\begin{equation}
f_R \approx \frac{\sum_i w_i\,P_i\cos^2\beta_i}{\sum_i w_i\,P_i},
\end{equation}
где $\beta_i$ --- угол нормали к плоскости $i$ с осью $c$, $w_i$ --- её доля телесного угла. На рис.~\ref{fig:t4_tc} это
проделано для расчётной текстуры: по отражениям до $2\theta = 80^\circ$ получается $f_R\approx0.60$, по отражениям до
140\textdegree{} --- 0.62, точное значение --- 0.61. Чем больше отражений и чем ближе они к максимумам текстуры, тем точнее.

\fig{t4_tc}{Обратная полюсная фигура по дифрактограмме. а) Плотности полюсов $P_{hkil}$ отражений $\alpha$-Zr при
съёмке с поверхности трубы (Cu $K\alpha$), упорядоченные по углу нормали с осью $c$ (в скобках); синие --- отражения до
$2\theta = 80^\circ$, серые --- до 140\textdegree. б) Стандартный треугольник, разбитый на ячейки Вороного вокруг
измеренных полюсов (красные точки); цвет --- $P$ ячейки. Оценка $f_R$ по формуле Кернса и точное значение.}

\section{Откуда берётся текстура}

\textbf{Деформационная текстура.} Когда зерно деформируется скольжением, его решётка поворачивается: скольжение
меняет форму зерна, а соседи заставляют его подстраиваться под общую деформацию, и разница компенсируется
поворотом. Ориентации «стекаются» к устойчивым положениям, зависящим от схемы деформации и систем скольжения. В
цирконии основное скольжение --- призматическое, оно не укорачивает зерно вдоль $c$. Поэтому оси $c$ уходят в
направление, вдоль которого металл \emph{сжимается} сильнее всего: тогда деформацию можно обеспечить без дорогого
пирамидального скольжения.

\textbf{Прокатка труб.} Трубы катают на станах холодной прокатки (пильгерных), уменьшая одновременно диаметр и толщину
стенки. Соотношение задают коэффициентом
\begin{equation}
Q = \frac{\ln(t_0/t)}{\ln(D_0/D)},
\end{equation}
где $t$ --- толщина стенки, $D$ --- средний диаметр. При $Q > 1$ сильнее утоняется стенка --- главное сжатие по радиусу,
и оси $c$ уходят к радиусу ($f_R$ растёт). При $Q < 1$ сильнее уменьшается диаметр --- главное сжатие по окружности,
и оси $c$ уходят к окружному направлению (рис.~\ref{fig:t5_pilger}). Оболочки катают с $Q > 1$, чтобы получить
радиальную текстуру --- как раз ради окружных гидридов.

\fig{t5_pilger}{Схема влияния коэффициента $Q$ при прокатке труб. а) $Q > 1$: в основном утоняется стенка, оси $c$
(красные отрезки, по два направления $\pm\chi$) собираются около радиуса. б) $Q < 1$: в основном уменьшается
диаметр, оси $c$ уходят к окружному направлению.}

\textbf{Отжиг.} При отжиге для снятия напряжений текстура почти не меняется. При рекристаллизации новые зёрна
наследуют распределение осей $c$, но поворачиваются вокруг $c$ примерно на 30\textdegree{} (рис.~\ref{fig:t2_polefigs}г).
Сильный рост зерна и особенно переход через $\beta$-область (нагрев выше 860--900\degC{} и охлаждение) текстуру
размывают: из одного зерна $\beta$ по соотношению Бюргерса возникает до 12 вариантов ориентации $\alpha$.

\section{Текстура и гидриды}

Гидрид ложится около базисной плоскости своего зерна (главы~\ref{ch:lattice} и \ref{ch:elastic}). Значит,
распределение нормалей к пластинкам гидридов внутри зёрен повторяет базисную полюсную фигуру, слегка размытую на
15\textdegree{} (плоскость $\{10\bar17\}$ отклонена от базиса).

\begin{zrbox}
Теперь всё складывается. В трубе с радиальной текстурой базисные плоскости большинства зёрен лежат примерно вдоль
окружности, поэтому без нагрузки гидриды \emph{окружные}. Окружное напряжение облегчает выпадение в зёрнах, у
которых ось $c$ ближе к окружному направлению (глава~\ref{ch:elastic}): в них пластинка встаёт по радиусу. Таких
зёрен немного --- им соответствует $f_T$. Поэтому переориентация зависит от текстуры. Дополнительный путь дают
межзёренные гидриды и цепочки пластинок, переходящие из зерна в зерно.
\end{zrbox}

\begin{warnbox}
Факторы Кернса --- средние $\cos^2$, и разные текстуры могут давать одинаковые $f$. Для переориентации важна не только
средняя, но и \emph{хвостовая} доля зёрен с осью $c$ почти по окружности: именно в них радиальные гидриды
зарождаются легче всего. Поэтому модель берёт полное распределение ориентаций (в простейшем случае --- наклон $\chi$
и разброс), а не три числа.
\end{warnbox}

\begin{exercise}
Оси $c$ всех зёрен лежат в плоскости радиус--окружность под углом $\pm\chi$ к радиусу. Найдите $f_R$ и $f_T$.
При каком $\chi$ они равны?
\end{exercise}

\begin{exercise}
Ориентация зерна задана углами Бунге $(\varphi_1, \Phi, \varphi_2) = (90^\circ, 30^\circ, 0)$; оси образца $X$ --- ось
трубы, $Y$ --- окружность, $Z$ --- радиус. Куда направлена ось $c$? Чему равны вклады этого зерна в $f_R$, $f_T$, $f_A$?
\end{exercise}

\begin{exercise}
Коэффициенты теплового расширения $\alpha$-Zr: $\alpha_a = 5.5\cdot10^{-6}$, $\alpha_c = 10.5\cdot10^{-6}$~К$^{-1}$. Оцените
коэффициенты расширения трубы по радиусу, окружности и оси при $f_R = 0.61$, $f_T = 0.35$, $f_A = 0.04$.
\end{exercise}
